\par\medskip
\SourceLeaf{078}\pdfbookmark[2]{59.}{article-059}
\begin{center}59.\end{center}
Die Regeln zur Bestimmung der Abstände von Ebenen, die durch die Sonne gehen, lassen sich offenbar auch auf die Abstände der Erde anwenden. Dabei kommen jedoch gewöhnlich nur die einfachsten Fälle vor. Es sei $R$ der Abstand der Erde von der Sonne, $L$ die heliozentrische Länge der Erde (die sich um $180^0$ von der geozentrischen Länge der Sonne unterscheidet) und $X$, $Y$, $Z$ die Abstände der Erde von drei Ebenen, die sich in der Sonne rechtwinklig schneiden. Wir betrachten nun die folgenden Fälle.

I. Ist die Ebene für die Abstände $Z$ die Ekliptik selbst und liegen die Pole der beiden anderen Ebenen, von denen die Erde die Abstände $X$, $Y$ hat, auf den Längen $N$ beziehungsweise $N+90^0$, so gilt
\[
X=R\cos(L-N),\qquad Y=R\sin(L-N),\qquad Z=0.
\]

II. Ist die Ebene für die Abstände $Z$ zum Äquator parallel und betragen die Rektaszensionen der Pole der beiden anderen Ebenen, von denen die Erde die Abstände $X$, $Y$ hat, $0$ beziehungsweise $90^0$, so erhält man, wenn $\epsilon$ die Schiefe der Ekliptik bezeichnet,
\[
X=R\cos L,\qquad Y=R\cos\epsilon\sin L,\qquad Z=R\sin\epsilon\sin L.
\]
Die namhaften Herausgeber der neuesten Sonnentafeln, \textsc{de Zach} und \textsc{de Lambre}, haben begonnen, die Breite der Sonne zu berücksichtigen. Diese Größe entsteht durch die Störungen der übrigen Planeten und des Mondes und kann kaum eine Bogensekunde erreichen. Bezeichnen wir die heliozentrische Breite der Erde, die stets denselben Betrag wie die Breite der Sonne und das entgegengesetzte Vorzeichen hat, mit $B$, so erhalten wir
\[
\begin{array}{@{}l|l@{}}
\text{in Fall I.}&\text{in Fall II.}\\[.4em]
X=R\cos B\cos(L-N)&X=R\cos B\cos L\\[.4em]
Y=R\cos B\sin(L-N)&Y=R\cos B\cos\epsilon\sin L-R\sin B\sin\epsilon\\[.4em]
Z=R\sin B&Z=R\cos B\sin\epsilon\sin L+R\sin B\cos\epsilon.
\end{array}
\]
Hier kann man $\cos B$ stets unbedenklich durch $1$ ersetzen und den in Teilen des Radius ausgedrückten Winkel $B$ für $\sin B$ einsetzen.

Die so gefundenen Koordinaten beziehen sich auf den \textit{Mittelpunkt} der Erde. Sind $\xi$, $\eta$, $\zeta$ die Abstände eines beliebigen Punktes auf der Erd\textit{oberfläche} von drei Ebenen durch den Erdmittelpunkt, die zu den durch die Sonne gelegten Ebenen parallel sind, so betragen die Abstände dieses Punktes von den \SourceLeaf{079}Ebenen durch die Sonne offenbar $X+\xi$, $Y+\eta$, $Z+\zeta$. Die Werte der Koordinaten $\xi$, $\eta$, $\zeta$ lassen sich in beiden Fällen leicht auf folgende Weise bestimmen. Es sei $\rho$ der Radius der Erdkugel (oder der Sinus der mittleren Horizontalparallaxe der Sonne) und $\lambda$ die Länge des Punktes auf der Himmelskugel, auf den die Gerade vom Erdmittelpunkt zu dem Punkt auf der Oberfläche projiziert wird. Seine Breite sei $\beta$, seine Rektaszension $\alpha$ und seine Deklination $\delta$. Dann gilt
\[
\begin{array}{@{}l|l@{}}
\text{in Fall I.}&\text{in Fall II.}\\[.4em]
\xi=\rho\cos\beta\cos(\lambda-N)&\xi=\rho\cos\delta\cos\alpha\\[.4em]
\eta=\rho\cos\beta\sin(\lambda-N)&\eta=\rho\cos\delta\sin\alpha\\[.4em]
\zeta=\rho\sin\beta&\zeta=\rho\sin\delta.
\end{array}
\]
Dieser Punkt auf der Himmelskugel fällt offenbar mit dem Zenit des Ortes auf der Erdoberfläche zusammen, sofern man die Erde als Kugel betrachtet. Seine Rektaszension stimmt daher mit der Rektaszension der Himmelsmitte, also mit der in Grad umgerechneten Sternzeit, überein und seine Deklination mit der Polhöhe. Wäre es der Mühe wert, die sphäroidische Gestalt der Erde zu berücksichtigen, so müsste man für $\delta$ die \textit{korrigierte} Polhöhe und für $\rho$ den wahren Abstand des Ortes vom Erdmittelpunkt ansetzen. Beide Größen lassen sich nach bekannten Regeln bestimmen. Aus $\alpha$ und $\delta$ ergeben sich die Länge und Breite $\lambda$ und $\beta$ nach den bekannten Regeln, die auch weiter unten angegeben werden. Im Übrigen ist klar, dass $\lambda$ mit der Länge des \textit{Nonagesimus} und $90^0-\beta$ mit dessen Höhe übereinstimmt.
